\section{Introduction}
A failed test is an observation, not a diagnosis. Suppose a program returns the wrong boundary value. An off-by-one error, an incorrect loop invariant, and a parsing mistake may all explain that outcome, yet suggest incompatible edits. Additional tests can distinguish these explanations, but only if their evidence is associated with the program version that produced it. Once the code changes, some explanations should survive, others should lose probability, and new failure mechanisms can appear. This motivates maintaining uncertainty over the cause of failure instead of treating execution feedback as a definitive instruction.

Language models already exploit execution feedback in several ways. Self-Debugging conditions refinement on program explanations and execution results~\citep{selfdebug}; REx allocates repair effort among candidates using an exploration--exploitation strategy~\citep{rex}; execution-feedback reinforcement learning changes the policy to improve iterative synthesis~\citep{rlef}. These approaches address important aspects of repair, but a transcript or a candidate-level success score is not itself an explicit calibrated posterior over latent explanations. Conversely, probabilistic inference over generated rollouts~\citep{roulette} represents uncertainty over a different object from the hidden failure state of a fixed candidate. We do not assume these alternatives cannot encode diagnostic information implicitly; we test whether exposing it as a predictive belief adds value.

We propose Particle-Belief Program Filtering (PBPF). Each candidate version owns a weighted collection of latent failure particles. The latent is piecewise static: ordered tests change its weights while source code is unchanged, whereas a patch creates a new child state linked only to its selected parent. A learned likelihood predicts five execution categories. Importance correction accounts for evidence-conditioned proposals at creation, and effective-sample-size-triggered resample--move controls particle degeneracy. This is an application of established sequential Monte Carlo principles~\citep{delmoral2006,gilks2001}, not a claim to invent Bayesian debugging.

The belief must also induce a coherent action. PBPF samples one particle and uses its learned soft prefix for a whole repair continuation. It does not decode a full program per diagnostic particle. Actor training evaluates a mixture of complete sequences; inference preserves a component within a sequence. This construction retains mutually incompatible hypotheses without averaging their latent representations or switching diagnoses token by token. Whether the resulting uncertainty is useful is an empirical question, not a consequence of the construction alone.

This prospective manuscript separates that question into prediction and intervention. Future-outcome prediction must improve over a development-selected deterministic memory before hidden-test repair is interpreted confirmatorily. All formal results remain pending. The proposed contributions are:
\begin{itemize}
\item A candidate-specific, piecewise-static diagnosis model with explicit root and child importance corrections, auditable normalizers, and isolated candidate lineages; an exact finite audit tests its inference behavior.
\item A coherent posterior-to-action interface that conditions one entire repair on one particle and trains a whole-sequence mixture; MAP, mean, remix, and ensemble controls test the role of preserved uncertainty.
\item A prospective, leakage-safe evaluation linking withheld-outcome prediction to equal-budget hidden repair, with source-cluster statistics, causal evidence controls, and locked model replications.
\end{itemize}
